% Generated from paper/full_draft. Edit the Markdown source or migration script.

\section{Method}
\label{sec:04-method:method}

\subsection{Problem Formulation}
\label{sec:04-method:problem-formulation}

In the retrieval-backed implementation evaluated here, let an evidence corpus be a set of chunks $D = \{d_i\}_{i=1}^{N}$ and a query stream $Q = \{q_t\}_{t=1}^{T}$. For each query $q_t$, a retrieval system returns an ordered context $C_t = [d_{t,1}, \ldots, d_{t,k}]$ that will be passed to a downstream generator. The objective is to preserve retrieval quality while controlling both retrieval cost and final context cost.

IntentRoute separates confidence-gated routing from final-context budgeting. For each query, the system estimates how much to rely on global dense retrieval, lexical BM25 recall, and cluster-local retrieval. A separately calibrated policy then budgets the resulting ranked evidence. Retrieval quality and final context tokens remain the primary mechanism-level outcomes; a frozen 300-query generation-and-judge experiment evaluates whether the same trade-off reaches answer-level behavior.

\subsection{Piecewise Relevance-Manifold Assumption}
\label{sec:04-method:piecewise-relevance-manifold-assumption}

The method is motivated by a bounded assumption:

\begin{quote}
In vertical-domain evidence retrieval, query-document relevance often follows a piecewise local structure induced by domain terminology, semantic neighborhoods, document organization, and user intent.
\end{quote}

This does not mean that corpus geometry is sufficient by itself. Instead, IntentRoute uses geometry as one routing signal among several. Dense retrieval remains a global recall floor, BM25 provides lexical anchors, and cluster-local retrieval provides local evidence patches.

\subsection{Multi-Route Retrieval Surface}
\label{sec:04-method:multi-route-retrieval-surface}

IntentRoute uses three retrieval routes.

\textbf{Dense route.} The dense route performs global dense retrieval over the selected corpus using cosine similarity in embedding space. It is the primary quality baseline and the main recall floor.

\textbf{BM25 route.} The BM25 route performs global lexical retrieval. It supports exact-match and terminology-sensitive queries and provides a lexical alternative when dense similarity is insufficient.

\textbf{Cluster-local route.} The cluster-local route performs dense retrieval within LinUCB-selected cluster arms. It allows the system to exploit local structure without searching the entire corpus as aggressively.

The routes are fused in the retrieval layer. The system can run in a full multi-route mode, a gated cost-aware mode, or a final context compaction mode. Full multi-route retrieval improves coverage but does not automatically reduce final context tokens. The final context policy described below is the mechanism that converts route confidence into retrieved-context token savings.

Figure~\ref{fig:system} summarizes this route-control architecture after the route surface is defined. It should be read as a controller diagram rather than a claim that LinUCB replaces dense or BM25 retrieval: dense and BM25 are global recall routes, LinUCB selects cluster-local arms, and route confidence is passed to the final context-budget controller.

\subsection{Routing Arm Construction}
\label{sec:04-method:routing-arm-construction}

Corpus chunk embeddings are clustered with KMeans or MiniBatchKMeans. This is a deliberate experimental choice. LinUCB requires a fixed number of arms, fixed arms improve reproducibility across seeds and scales, and KMeans is fast enough for large-scale LoTTE experiments. The same arm count is used across LoTTE scales to keep the LinUCB state space comparable, even though larger corpora therefore contain more chunks per arm. We use 32 routing arms as the main reproducible operating point. A sensitivity study over $K \in \{8,16,32,64,128\}$ shows that full multi-route fused quality is stable across this range, while retrieval-stage gated routing is more sensitive because smaller $K$ changes route granularity and larger $K$ spreads feedback across more arms. We therefore treat 32 as an engineering choice, not a manifold-derived optimum.

The paper does not claim that KMeans is the best clustering algorithm for retrieval. HDBSCAN, graph clusters, or learned routing structures may be better in some deployments, but dynamic arm counts complicate the current LinUCB setup. Here, each cluster arm represents a local region of the retrieval surface. The cluster route searches inside selected arms, while dense and BM25 routes can rescue cases where the cluster route misses relevant evidence.

\subsection{LinUCB Route Policy}
\label{sec:04-method:linucb-route-policy}

For each query $q_t$, IntentRoute computes a context feature vector $x_t \in \mathbb{R}^{p}$. Features include query embedding projections, route confidence signals, and local geometry signals. For each arm $a \in \mathcal{A}$, LinUCB maintains a linear value model:

\[
\begin{aligned}
\hat{\theta}_a &= A_a^{-1} b_a, \\
s_t(a) &= \hat{\theta}_a^\top x_t
        + \alpha \sqrt{x_t^\top A_a^{-1} x_t}.
\end{aligned}
\]

The first term estimates the arm value under the current query context. The second term is an exploration bonus. The parameter $\alpha$ controls the exploration-exploitation balance: larger values encourage the policy to explore arms whose value estimate is uncertain.

The policy selects the top candidate arms by score. The selected arms define the cluster-local retrieval path and provide confidence signals for later context compaction.

\subsection{Feature Groups and Route Confidence}
\label{sec:04-method:feature-groups-and-route-confidence}

The LinUCB context vector is not intended to introduce a new representation model. It collects signals that decide whether local routing is reliable for the current query. The following summary lists the feature groups used by the controller.

The feature groups are:

\begin{itemize}
\item query representation: normalized query embedding and PCA/context projection, used to place the query on the same local surface as corpus arms;
\item dense confidence: dense score concentration, top-rank margin, and fallback availability, used to estimate whether global dense retrieval is already reliable;
\item lexical confidence: BM25 candidate availability and lexical-match strength, used to protect terminology-heavy and exact-match queries;
\item route agreement: overlap among dense, BM25, and cluster-local candidates, used to detect when independent routes support the same evidence region;
\item local geometry: nearest centroid similarity, selected arm identity, and semantic drift, used to estimate whether the query lies close to selected local arms;
\item feedback state: selected-arm value, pull-count maturity, and recent route reward, used to estimate whether LinUCB has enough evidence to trust the local route;
\item budget state: route confidence tier and fallback status, used to decide whether final context compaction is allowed.
\end{itemize}

Route confidence is computed from the selected-arm value estimate, the top-versus-rest arm margin, and an arm-maturity term based on feedback count. Semantic drift is defined as one minus the nearest selected-centroid similarity. Low confidence or high drift keeps the dense fallback active; high confidence and low drift allow the final-context policy to compact evidence.

\subsection{Trust-Weighted Feedback}
\label{sec:04-method:trust-weighted-feedback}

IntentRoute models feedback as a noisy signal rather than a perfect oracle. In the trust-weighted mode, each simulated user feedback event is assigned a trust weight. Higher-trust feedback contributes more to the arm update and local feedback memory. Lower-trust feedback has a weaker effect.

Conceptually:

\[
\begin{aligned}
r_t &= g_t - \lambda c_t, \\
\tilde{r}_t &= \tau_t r_t, \\
A_a &\leftarrow A_a + \tau_t x_t x_t^\top, \\
b_a &\leftarrow b_a + \tilde{r}_t x_t.
\end{aligned}
\]

Here, $g_t$ is the retrieval-quality signal, $c_t$ is the cost penalty, $\lambda$ controls the quality-cost trade-off, $\tau_t$ is the feedback trust weight, and $\tilde{r}_t$ is the weighted reward applied to the selected arm.

The current experiments do not claim that real human feedback was collected. They show that under controlled simulated feedback, the route policy can self-improve. The strongest feedback evidence is visible in policy metrics such as last true reward and selected-cluster hit rate, especially when final retrieval quality is already protected by dense and BM25 fallback.

\subsection{Route-Level Credit Assignment}
\label{sec:04-method:route-level-credit-assignment}

Early experiments used final fused retrieval success as the reward signal. This can over-credit the selected cluster arm when dense or BM25 rescue the final ranking. IntentRoute therefore uses a stricter \nolinkurl{cluster_only} reward attribution mode in the main policy-learning experiments. The LinUCB arm is updated using the quality of its own cluster-local route.

This distinction is important. The final fused ranking measures the system outcome, while cluster-only reward measures whether LinUCB is learning a better route. Dense and BM25 rescue paths protect final quality; cluster-only credit assignment tests whether the adaptive component itself is improving.

\subsection{Prequential Adaptation Protocol}
\label{sec:04-method:prequential-adaptation-protocol}

The evaluation uses a no-leakage prequential protocol. For each query $q_t$, the current policy state is frozen before retrieval. The system ranks candidates, constructs the final context, and is evaluated against ground-truth evidence. Only after this evaluation is the ground-truth label converted into simulated feedback and used to update the LinUCB state for later queries.

This means feedback for $q_t$ cannot improve the ranking of $q_t$ itself. Earlier feedback can influence later queries, but future query feedback is not available to the current policy. The protocol should therefore be described as simulated test-time adaptation, not as offline IID held-out generalization.

Several repeated-feedback experiments use multiple prequential epochs over the same query stream to simulate repeated interactions. Each epoch preserves the same discipline: rank and evaluate the current query before applying feedback from that query. Multi-epoch results are therefore controlled repeated interaction studies. They should not be over-interpreted as single-pass generalization results.

\subsection{Confidence-Based Final Context Compaction}
\label{sec:04-method:confidence-based-final-context-compaction}

Preliminary token-cost analysis showed that reducing source candidates does not automatically reduce final context tokens. IntentRoute therefore combines confidence-gated routing with an explicit calibrated budget policy, the central quality-cost control interface.

Let $R_t$ be the ranking produced by the confidence-gated route surface. A calibration policy $\pi_\phi$ selects a token ratio $r \in (0,1]$ and minimum prefix $m$. The final context is the longest ranked prefix of at least $m$ chunks satisfying

\[
\mathrm{Tokens}(C_t) \le r\,\mathrm{Tokens}(R_t[:10]).
\]

The policy parameters $\phi$ are selected on calibration queries subject to a retrieval-quality eligibility gate and then frozen before held-out test evaluation. Geometry and feedback affect route construction, while the stronger token saving arises when the separately calibrated length budget acts on the routed ranking. The implementation does not learn a direct per-query mapping from confidence to token ratio.

The conservative \nolinkurl{confidence_topk} policy works as follows:

\begin{enumerate}
\item If route confidence is low or semantic drift is high, keep dense fallback and the normal top-10 context.
\item If LinUCB confidence is high, reduce final context size from the default top-10 to $k=8$.
\item If confidence is mid-level, keep $k=10$ as a safety tier rather than calling it true compression.
\item Report the actual retrieved context token count, not only candidate counts.
\end{enumerate}

For the conservative policy, the effective compaction rate is therefore the high-confidence compression rate, not the broader rate of queries that enter a non-fallback route. Hybrid-lite can reduce dense influence in fusion while retaining dense candidates as a safety net; it should not be described as reducing dense computation unless the global dense route is actually skipped.

The main token-quality result uses frozen calibration/test budget policies that select global ratio/minimum-prefix parameters on calibration queries and evaluate them unchanged on held-out test queries. The conservative confidence-based policy remains a stable baseline: it reduces final context tokens by about 4.7-5.3\% across LoTTE 100k, 200k, 400k, and 638k while preserving dense-level query hit.

On LoTTE, semantic drift rarely exceeds the configured fallback threshold, so the reported compression behavior is primarily confidence-driven. In more heterogeneous query distributions, drift-based fallback may become more active.

\subsection{Feedback-Triggered Recovery}
\label{sec:04-method:feedback-triggered-recovery}

IntentRoute can also use feedback as a recovery signal after a compressed context fails. In this mode, feedback does not insert the missing ground-truth chunk into the current answer. Instead, it updates arm-level route value and marks the associated local region as risky for future routing or optional post-feedback retry.

The safe recovery behavior is conservative:

\begin{enumerate}
\item If feedback indicates that a compressed context missed evidence, identify the evidence-bearing arm or risky query-arm relation.
\item Prefer a less aggressive final-context budget or full-context fallback for that local region.
\item Avoid unconditional global arm boosting, because a feedback-positive arm for one query can transfer poorly to unrelated queries.
\end{enumerate}

This recovery mechanism is deployment-facing. It represents how a system can repair tail failures after user or evaluator feedback, not how the first-pass ranking is produced before feedback is observed.

\subsection{Algorithm Sketch}
\label{sec:04-method:algorithm-sketch}

Input: query $q_t$, corpus $D$, route artifacts, and the current LinUCB state.

\begin{enumerate}
\item Embed $q_t$ and compute context features $x_t$.
\item Score cluster arms with LinUCB using $s_t(a)$.
\item Retrieve candidates from the global dense route, the global BM25 route, and dense search within selected cluster arms.
\item Fuse route rankings.
\item Apply confidence-based final context policy.
\item Evaluate retrieval quality for $q_t$.
\item Only after evaluation, convert the ground-truth label into simulated feedback and update LinUCB for later queries.
\item If the interaction is a post-feedback retry or a later query in a risky local region, use the updated arm state to choose a safer context budget or fallback path.
\end{enumerate}

Output: final retrieved context $C_t$ and updated policy state.

\subsection{Reproducibility Parameters}
\label{sec:04-method:reproducibility-parameters}

The following summary lists the main implementation parameters used in the reported cost-aware LinUCB experiments. Scale-specific cache paths and dataset sizes are reported in the experiment artifacts; these parameters define the controller behavior.

The main controller parameters are:

\begin{itemize}
\item KMeans/MiniBatchKMeans arms: 32 fixed LinUCB arms;
\item candidate arms per query: 3 cluster-local routes;
\item context projection dimension: 64;
\item LinUCB exploration: $\alpha=1.0$, decay 0.01, minimum 0.3;
\item prequential epochs: 3 unless otherwise stated;
\item feedback trust: default $\tau=0.75$ for noisy feedback updates;
\item full-route candidate depths: dense/BM25/cluster = 100/100/100;
\item lite-route depths: dense/BM25 = 20/20;
\item dense safety floor: 5 full-route chunks and 2 lite-route chunks;
\item route fusion: weighted reciprocal-rank fusion with $k=60$;
\item full-route weights: dense 2.0, BM25 0.8, cluster 0.8;
\item lite-route weights: dense 0.8, BM25 0.5, cluster 2.0;
\item confidence thresholds: high 0.65 and mid 0.35;
\item drift threshold: 1.0;
\item token-budget grid: $r \in \{0.85,0.88,0.90,0.92,0.95,0.98\}$ and $m \in \{4,\ldots,8\}$.
\end{itemize}

The notation \nolinkurl{token_budget_r0.85_m4} means that each query keeps a safe prefix of at least four chunks and then admits additional chunks only while the final context remains within 85\% of the original dense top-10 token budget. The policy is chosen on calibration queries and then frozen before held-out test evaluation.
